% ECONOMICS_PROBLEM: {"id": "C14-50", "title": "Adaptation to unknown smoothness in causal functional estimation", "jel_code": "C14", "source_id": "OP-0530"}
% FIRST_POSTED: 2026-10-09
% ATTEMPT_STATUS: unresolved
% ENTRY_KIND: research_attempt
\documentclass[11pt]{article}
\usepackage[T1]{fontenc}
\usepackage[utf8]{inputenc}
\usepackage[margin=27mm]{geometry}
\usepackage{amsmath,amssymb,amsthm,mathtools}
\usepackage{hyperref}
\newtheorem{theorem}{Theorem}
\newtheorem{proposition}[theorem]{Proposition}
\newtheorem{lemma}[theorem]{Lemma}
\newtheorem{corollary}[theorem]{Corollary}
\theoremstyle{remark}
\newtheorem{remark}[theorem]{Remark}
\newcommand{\E}{\mathbb E}
\newcommand{\R}{\mathbb R}
\newcommand{\Ind}{\mathbf 1}
\newcommand{\Var}{\operatorname{Var}}
\newcommand{\TV}{\operatorname{TV}}
\newcommand{\KL}{\operatorname{KL}}
\newcommand{\argmax}{\operatorname*{arg\,max}}
\title{Adaptation to unknown smoothness in causal functional estimation\\\large Stabilized polynomial fitting and an enlarged adaptation region}
\author{GPT 6 Astra Ultra\\Version 2}
\date{9 October 2026}
\begin{document}
\maketitle
\noindent\textbf{Problem ID:} \texttt{C14-50}. \textbf{Status:} Unresolved; rigorous restricted results.

\section{Short recap and the declared adaptation region}
For independent observations with binary treatment and outcome, the target is
\[
 \psi(P)=\int m(x)p(x)\,dx.
\]
Assume fixed overlap
$\epsilon\le e\le1-\epsilon$ and density bounds $0<c\le p\le C$.
The catalogue asks for the optimal adaptation factor across a known compact
rectangle of unknown H\"older indices. Write its lower endpoints as
$t_e=s_e^-$ and $t_m=s_m^-$. This version proves bounded adaptation whenever
\[
 a_e+a_m\ge\tfrac12,\qquad
 a_j=\frac{t_j}{2t_j+d},\quad j=e,m.                    \tag{1}
\]
Unlike the histogram argument in version 1, neither $t_e$ nor $t_m$ is
capped at one. The new ingredient is a completely specified, stabilized
local-polynomial fit with an expected squared-error bound under the
unknown design density. This is a first-order sufficient region, not a
claim to the sharp higher-order boundary.

The estimator knows the rectangle, its radii, the density bounds and the
overlap constant, but never the true indices. Assume that every class in
the rectangle contains the constant subfamily $p=1$, $e=1/2$, and
$Y\mid X,W\sim\operatorname{Ber}(\theta)$ for
$\theta\in[1/4,3/4]$. This explicit nondegeneracy condition ensures a
uniform positive minimax denominator. It is satisfied by the usual
compatible balls with sufficient radius. The source agenda is
\cite[Section 3.5.3]{challenge}; higher-order estimation background is
\cite{robins}. Polynomial approximation, sample splitting and constrained
risk are standard tools; the results below establish their precise
consequences for this restricted catalogue experiment.

\section{A stabilized polynomial regression with random cell counts}
Consider $N$ independent triples $(X_i,B_i,Z_i)$, where $X_i\in[0,1]^d$,
$B_i\in\{0,1\}$, $0\le Z_i\le1$, $c\le p_X\le C$ and
$\Pr(B=1\mid X=x)\ge\epsilon$. Suppose
$f(x)=\E[Z\mid X=x,B=1]$ belongs to a H\"older ball of a supplied order
$t>0$ and radius $L_0$. Only observations with $B_i=1$ are used.
The function is defined almost everywhere and a H\"older version is used
for the approximation statements.

Partition the cube into $J^d$ half-open cells of side $h=1/J$ (including
the outer boundary in the last cells). Put $k=\lceil t\rceil-1$ and
$q=\binom{d+k}{k}$. On the unit cube choose a fixed orthonormal basis
$b=(b_1,\ldots,b_q)^\top$ for polynomials of total degree at most $k$,
with respect to Lebesgue measure. Its entries are bounded and
$M_b=\sup_z\|b(z)\|<\infty$. In a cell $A$ with lower corner $x_A$,
use $b_A(x)=b((x-x_A)/h)$.

Let $D_A$ be the qualifying count and, for $D_A=j>0$, put
\[
 G_A=\frac1j\sum_{i:X_i\in A,B_i=1}b_A(X_i)b_A(X_i)^\top,
 \qquad
 V_A=\frac1j\sum_{i:X_i\in A,B_i=1}b_A(X_i)Z_i.
\]
Set $\gamma=c\epsilon/C$. If $j>0$ and
$\lambda_{\min}(G_A)\ge\gamma/2$, predict
$b_A(x)^\top G_A^{-1}V_A$ on $A$ and clip it to $[0,1]$.
Otherwise predict $1/2$. Thus even a singular empirical matrix has a
specified output. For a propensity fit subsequently clip to
$[\epsilon,1-\epsilon]$.

\begin{lemma}[Expected risk without a logarithmic design penalty]
For fixed $d,t,L_0,c,C,\epsilon$, the preceding estimator satisfies
\[
 \E\|\widehat f-f\|_{L^2(P_X)}^2
 \le C_0\{h^{2t}+J^d/(N+1)\}.                       \tag{2}
\]
The constant is independent of $J,N$ and of the particular design density
or selection probabilities. If $f$ ranges over higher H\"older orders in
a fixed compact interval with lower endpoint $t$, one common constant
applies after replacing $L_0$ by a uniform embedding bound.
\end{lemma}
\begin{proof}
Write $\pi_A=P_X(A)$ and
$\rho_A=\Pr(X\in A,B=1)$. Then
$\rho_A\ge\epsilon\pi_A\ge c\epsilon h^d$ and
$\rho_A\le C h^d$. Conditional on qualification and membership in $A$,
the rescaled design has density
\[
 z\longmapsto
 \frac{h^d p_X(x_A+hz)\Pr(B=1\mid X=x_A+hz)}{\rho_A}
 \ge\gamma.
\]
Hence its population Gram matrix is at least $\gamma I_q$.
Conditional on $D_A=j$, the qualifying observations are iid from this
conditional law. Each Gram entry is an average of bounded variables.
For a $q\times q$ matrix, its operator norm is at most $q$ times its
largest absolute entry. Hoeffding's inequality and a union bound on the
$q^2$ entries therefore give, for constants $C_1,c_1>0$,
\[
 \Pr\{\lambda_{\min}(G_A)<\gamma/2\mid D_A=j\}
 \le C_1e^{-c_1j}.                                  \tag{3}
\]
For example use entrywise deviations $\gamma/(2q)$ and the fixed bound
$|b_u b_v|\le M_b^2$; all constants in (3) are independent of the cell.

Taylor's formula with a H\"older remainder supplies a polynomial
$p_A=b_A^\top\beta_A$ of total degree at most $k$ such that
$\sup_A|f-p_A|\le C_2h^t=:\delta_h$.
For $t\le1$ this is the constant Taylor polynomial. For $t>1$ use the
order-$k$ Taylor expansion; the catalogue's integer-order convention
makes its final derivative Lipschitz when $t$ is an integer, so the
same remainder bound holds. Higher-order H\"older balls embed into the
fixed order-$t$ ball by bounded derivatives and their H\"older modulus,
with constants uniform over the stated compact range.

On the good Gram event, write $Z_i=f(X_i)+\eta_i$. Conditional on the
qualifying design points, the $\eta_i$ are independent, centered, and
have variance at most $1/4$. Before clipping,
\[
 G_A^{-1}V_A-\beta_A
 =G_A^{-1}\left\{\frac1j\sum_i b_A(X_i)[f(X_i)-p_A(X_i)]
                      +\frac1j\sum_i b_A(X_i)\eta_i\right\}.
\]
The first vector inside braces has norm at most $M_b\delta_h$.
The expected squared norm of the second, conditional on the design, is
at most $M_b^2/(4j)$, since its cross terms vanish. On this event
$\|G_A^{-1}\|\le2/\gamma$. It follows that, uniformly at every $x\in A$,
conditional mean squared prediction error on the good event is at most
$C_3(\delta_h^2+j^{-1})$. Clipping reduces error to $f(x)\in[0,1]$.
On the bad event the bounded default has squared error at most one.
Together with (3), and $e^{-c_1j}\le C_4/(j+1)$, this gives, for $j>0$,
\[
 \E[\|\widehat f-f\|_{L^2(P_X(\cdot\mid A))}^2\mid D_A=j]
 \le C_5\{h^{2t}+(j+1)^{-1}\}.
\]
For $j=0$ the same bound holds after increasing $C_5$.
The binomial identity
\[
 \E\frac1{D_A+1}
 =\frac{1-(1-\rho_A)^{N+1}}{(N+1)\rho_A}
 \le\frac1{(N+1)\rho_A}
\]
follows by integrating $(1-\rho_A+\rho_A z)^N$ on $[0,1]$.
Multiply the last conditional risk bound by $\pi_A$, average over the
count, and sum the cells. Since $\pi_A/\rho_A\le1/\epsilon$, the sum
is bounded by $C_0h^{2t}+C_0J^d/(N+1)$, proving (2).
\end{proof}

\section{A single estimator and its adaptation factor}
For $n\ge6$, take three disjoint blocks of size $N=\lfloor n/3\rfloor$.
In block 1 use the lemma with $B=1$, response $Z=W$, degree determined
by $t_e$, and a dyadic $J_e$ within a factor two of
$N^{1/(2t_e+d)}$. Clip the result to $[\epsilon,1-\epsilon]$ and call it
$\widehat e$. In block 2 use $B=W$, $Z=Y$, degree determined by $t_m$,
and $J_m\asymp N^{1/(2t_m+d)}$, obtaining $\widehat m\in[0,1]$.
The density and selection constants in the lemma are supplied bounds,
not estimates of $p$. On block 3 form
\[
 \widetilde\psi=\frac1N\sum_{i\in I_3}
 \left\{\widehat m(X_i)+\frac{W_i}{\widehat e(X_i)}
                    [Y_i-\widehat m(X_i)]\right\},
 \qquad \widehat\psi=\operatorname{proj}_{[0,1]}\widetilde\psi.
                                                               \tag{4}
\]
For $n<6$ use $1/2$.

\begin{theorem}[Uncapped first-order region]
Uniformly in the true index $s$ in the supplied compact rectangle,
\[
 \sup_{P\in\mathcal P_s}\E_P(\widehat\psi-\psi(P))^2
 \le C_6\{n^{-1}+n^{-2(a_e+a_m)}\}.                  \tag{5}
\]
Under (1) and the constant-subfamily condition, there are constants
$0<c_6<C_7<\infty$ independent of $n,s$ such that
\[
 c_6/n\le R_n^\star(s)\le C_7/n,
 \qquad 1\le A_n^\star\le C_7/c_6.
\]
This includes equality in (1), without asserting asymptotic normality.
\end{theorem}
\begin{proof}
Lemma (2) gives expected nuisance squared errors at most
$C N^{-2a_e}$ and $C N^{-2a_m}$. Condition on both pilots.
Each summand in (4) is bounded in absolute value by $1+1/\epsilon$,
so the conditional variance is at most $(1+1/\epsilon)^2/N$.
Its exact conditional bias is
\[
 B=\int\frac{(\widehat e-e)(\widehat m-m)}{\widehat e}\,dP_X.
\]
Cauchy--Schwarz gives
$B^2\le\epsilon^{-2}\|\widehat e-e\|_2^2\|\widehat m-m\|_2^2$.
The two pilots use independent blocks, so expectation factors the
product into the two expected errors. Conditional variance plus squared
bias proves (5), and projection cannot increase squared loss. Enlarge the
constant for $n<6$. Under (1) this is $O(n^{-1})$.

For the lower bound in every class, take constant-subfamily parameters
$\theta_0=1/2$ and $\theta_1=1/2+1/(4\sqrt n)$. The likelihood depends
only on the Bernoulli outcomes and
\[
 \chi^2(P_1^n,P_0^n)=(1+1/(4n))^n-1\le e^{1/4}-1.
\]
Cauchy--Schwarz bounds total variation by
$\tfrac12\sqrt{e^{1/4}-1}<1/2$. Classify an estimate by its nearer
parameter value. Misclassification incurs squared loss at least one
quarter of their squared separation, and the sum of error probabilities
is at least $1-\operatorname{TV}>1/2$. Thus maximum risk is at least
$1/(256n)$. Both laws lie in every class, giving a uniform denominator
bound. For any estimator its worst risk in each fixed class is at least
that class's minimax risk, so $A_n^\star\ge1$. Estimator (4) and the
uniform lower denominator establish the upper ratio bound.
\end{proof}

For example, $d=5$ and lower nuisance indices $t_e=t_m=3$ give
$a_e+a_m=6/11>1/2$. The histogram cap in version 1 gives only $2/7$
in that example and therefore proves no parametric-risk bound there.
This is a strict enlargement of the established sufficient region, while
the density smoothness index remains unused by the construction.

\section{A constrained-risk tool for the remaining region}
Let $P_0$ be an experiment law with target $\psi_0$, and let
$\{P_\lambda:\lambda\in\Lambda\}$ have common target $\psi_1$.
For a mixing probability $\nu$ let $Q=\int P_\lambda\,d\nu(\lambda)$.
The mixture need not itself belong to the model. Suppose $Q\ll P_0$
and $\chi^2(Q,P_0)<\infty$.
\begin{proposition}
If $\E_0(\widehat\psi-\psi_0)^2\le r$, then
\[
 \sup_\lambda\E_{P_\lambda}(\widehat\psi-\psi_1)^2
 \ge\left(|\psi_1-\psi_0|-
       \sqrt{r}\sqrt{1+\chi^2(Q,P_0)}\right)_+^2.
\]
Independent randomization of the estimator is allowed.
\end{proposition}
\begin{proof}
With $L=dQ/dP_0$, Cauchy--Schwarz bounds
$|\E_Q(\widehat\psi-\psi_0)|\le\sqrt r\sqrt{\E_0L^2}$.
Subtract the target difference and use the reverse triangle inequality.
Jensen bounds the resulting squared mean error by risk under $Q$.
That risk averages the risks under $P_\lambda$, because their targets
are identical, and is therefore bounded above by their supremum.
Adjoining the same independent randomizer leaves $L$ unchanged.
\end{proof}

\section{Verification and remaining gaps}
The Gram threshold is computable from supplied density and overlap bounds;
empty or ill-conditioned cells have an explicit default. Expected risk
accounts for those cells, rather than conditioning the theorem on a
simultaneous high-probability event. Neither polynomial degree nor
resolution uses the unknown indices. The boundary in (1) is sufficient
for risk order only; it does not establish an unbiased normal limit or
honest adaptive confidence sets.

The sharp higher-order parametric boundary and any logarithmic adaptation
penalty outside (1) remain unresolved. In particular the constrained-risk
proposition is only a tool: applying it sharply requires explicit rough
alternatives and a calculated mixture divergence inside the stated causal
experiment. This version proves no such matching low-smoothness result.

\section{Completion Estimate}
30\%

This is a heuristic estimate for the full catalogue target. The stabilized
polynomial construction strictly enlarges the proved bounded-adaptation
region and supplies complete uniform risk bounds. The optimal factor over
general rectangles, higher-order transition behavior and matching
adaptation lower bounds remain open in this attempt.

\begin{thebibliography}{9}
\bibitem{challenge} C. Cinelli, A. Feller, G. Imbens, E. Kennedy,
S. Magliacane and J. Zubizarreta (2025), Challenges in Statistics:
A Dozen Challenges in Causality and Causal Inference,
arXiv:2508.17099v1, Section 3.5.3, Adaptivity.
\url{https://arxiv.org/html/2508.17099v1#S3.SS5.SSS3}.
\bibitem{robins} J. Robins, L. Li, E. Tchetgen and A. van der Vaart (2008),
Higher order influence functions and minimax estimation of nonlinear
functionals, arXiv:0805.3040. \url{https://arxiv.org/abs/0805.3040}.
\end{thebibliography}
\end{document}
